\section{Related Work}

\textbf{Red-Teaming and Jailbreaking LLMs.}
Early attempts at red-teaming and understanding LLM vulnerabilities have focused on hand-crafting prompts and analyzing model responses~\citep{hh-rlhf-preference, hh-rlhf, openai2024gpt4}. However, manual methods are limited in scope and efficiency due to the prohibitive cost of human annotations. Thus, automated red-teaming and jailbreaking methods are developed for more scalable audit of model vulnerabilities. One genre of methods involves gradient optimization that requires back-propagating through model parameters \citep{gcg,gbda,coldattack,schwinn2024soft}. However, they are computationally expensive, cannot be applied to closed models, and often result in gibberish texts. There are also inference-based approaches (most related to our work) which generate jailbreaking prompts directly or through iterative edits~\citep{pair,liu2023autodan,lapid2023open,li2024deepinception,red-teaming-lm-w-lm, casper2023explore,tap,yu2023gptfuzzer,promptpacker,yuan2023gpt4,zeng2024johnny,Deng_2024}. Other jailbreaking works study attacks during decoding time (\eg decoding configurations \citep{huang2023catastrophic}, logit manipulation \citep{zhao2024weaktostrong}), in other modalities \cite{shayegani2024jailbreak}, under multilingual settings \citep{deng2024multilingual,yong2024lowresource,qiu2023latent}, or in programming mode \cite{kang2023exploiting}. However, most jailbreak methods rarely result in large-scale training resources for model safety enhancement due to their limited coverage of attack and risk types and slow speed. \method differs from previous works by efficiently composing various adversarial attacks utilizing real-world jailbreak tactics mined from in-the-wild user-chatbot interactions. \method allows scalable synthetic safety training data generation in addition to simply showing its attack efficacy.

\textbf{Safety Evaluation and Enhancement of LLMs.}
Many red-teaming efforts on LLMs have been formalized as benchmarks for evaluating model vulnerabilities---these typically are composed of harmful prompts that models should refuse~\citep{adversariallyaligned,wei2023jailbroken,wang2023donotanswer,sun2024trustllm,mazeika2024harmbench,geiping2024coercing,wang2024decodingtrust,chao2024jailbreakbench}. Meanwhile, to mitigate the potential byproducts of safety training, other benchmarks measures exaggerated safety behavior on benign queries \citep{xstest,cui2024orbench}. While LLM safety evaluation has been an active area of research, studies and resources for safety training have been limited~\citep{hh-rlhf,dai2023safe,yang2023alignment}. Most related to our work in this space are \safellama \citep{bianchi2023safety} and BeaverTails~\citep{beavertails}, which primarily focus on vanilla harmful queries by releasing small-scale safety training datasets and large-scale pairwise preference datasets, respectively. \method distinguishes from these works by releasing higher quality (shown by our training ablation experiments) and larger scale sequential instruction-tuning data comprised of both vanilla and adversarial queries. Finally, synthetic data has been used for LLM safety \cite{casper2023explore,perez2022discovering,hubinger2024sleeper}. Most relevant to our work is Rainbow Teaming~\cite{rainbowteaming}, which uses synthetic data to populate a grid of attack spaces based on the attack style and risk category. Our work differs in automatically mining human-devised jailbreak tactics rather than manually defining attack styles \cite{rainbowteaming}, creating a large-scale open safety training resource that supports extensive safety training experiments.
